\documentclass[fleqn]{ctexart}
\usepackage{graphicx}
% 数学支持（提供 \mathbb、\dfrac 等）
\usepackage{amsmath,amssymb}
% 页面与版式设置：更小的页边距与顶部空白
\usepackage[a4paper,top=15mm,bottom=20mm,left=12mm,right=12mm]{geometry}
% 自动换行与字偶距优化
\usepackage{microtype}
\microtypesetup{tracking=true,protrusion=true,final}
% 在需要时放宽断行以避免溢出（数值可按需调整）
\setlength{\emergencystretch}{2em}
% 列表控制，便于让(1)后直接跟内容
\usepackage{enumitem}
% 使章节标题左对齐（同时保留 ctex 默认样式）
\ctexset{section={format+=\raggedright}}
%1.3倍行距
\renewcommand{\baselinestretch}{1.3}
% 数学左对齐缩进为0
\setlength{\mathindent}{0pt}
% 增大矩阵的行距
\renewcommand{\arraystretch}{1.3}
\pagestyle{empty}
\begin{document}
\section*{题目1}

分步骤解读代码：
\begin{enumerate}[label=(\arabic*)]
	\item 父进程在 \texttt{fork()} 之前先输出一次 \texttt{counter=2}
	\item \texttt{fork()} 后父子进程各自拥有独立的地址空间：两边的 \texttt{counter} 初值都为 2
	\item 子进程收到 \texttt{SIGUSR1} 后转去执行信号处理函数 \texttt{handler1}：
				\begin{itemize}
					\item 子进程把自己的 \texttt{counter} 从 2 减到 1，并输出 \texttt{1}
					\item 随后 \texttt{exit(0)} 使子进程直接退出
				\end{itemize}
	\item 父进程 \texttt{kill(pid, SIGUSR1)} 后立刻 \texttt{waitpid} 等待子进程结束；由于子进程在处理函数里 \texttt{exit(0)}，父进程的 \texttt{waitpid} 会返回
	\item 父进程返回后将自己的 \texttt{counter} 从 2 加到 3，最后输出 \texttt{3}
\end{enumerate}

\textbf{因此输出顺序是：}\ \texttt{2}（父）$\to$\texttt{1}（子）$\to$\texttt{3}（父），最终输出为：213

\section*{题目2}
实现思路：\texttt{sleep(n)} 的返回值是“剩余秒数”
\begin{itemize}
	\item 正常睡满：返回 0
	\item 被 Ctrl+C（\texttt{SIGINT}）打断：\texttt{sleep} 提前返回剩余秒数
\end{itemize}
为了不让 Ctrl+C 触发默认的“直接终止进程”，设置一个空的 \texttt{SIGINT} 处理函数即可

\begin{verbatim}
void sigint_handler(int sig) {
    (void)sig;
}

int main(int argc, char **argv) {
    int secs = atoi(argv[1]);
    signal(SIGINT, sigint_handler);

    unsigned int left = sleep((unsigned int)secs);
    printf("Slept for %u of %d secs.\n", (unsigned int)secs - left, secs);
    return 0;
}
\end{verbatim}

\section*{题目3}

\subsection*{(1) FCFS}
调度顺序：$1\to 2\to 3\to 4\to 5$

完成时刻：$C_1=8, C_2=11, C_3=13, C_4=23, C_5=24$

\begin{itemize}
	\item 周转：$T_1=8,\ T_2=11,\ T_3=13,\ T_4=21,\ T_5=22$，平均 $\bar T=\frac{75}{5}=15$
	\item 响应：$R_1=0,\ R_2=8,\ R_3=11,\ R_4=11,\ R_5=21$，平均 $\bar R=\frac{51}{5}=10.2$
\end{itemize}

\subsection*{(2) RR}
调度序列：

\medskip
$
\begin{aligned}
&1,2,3,1,2,4,5,3,\\
&1,2,4,1,4,1,4,1,\\
&4,1,4,1,4,4,4,4
\end{aligned}
$

\medskip
完成时刻：$C_5=7, C_3=8, C_2=10, C_1=20, C_4=24$

\begin{itemize}
	\item 周转：$T_1=20,\ T_2=10,\ T_3=8,\ T_4=22,\ T_5=5$，平均 $\bar T=\frac{65}{5}=13$
	\item 响应：
				$R_1=0,\ R_2=1,\ R_3=2,\ R_4=3,\ R_5=4$，平均 $\bar R=\frac{10}{5}=2$
\end{itemize}

\subsection*{(3) SJF}
调度顺序：
$3\to 5\to 2\to 1\to 4$

完成时刻：$C_3=2, C_5=3, C_2=6, C_1=14, C_4=24$

\begin{itemize}
	\item 周转：$T_1=14,\ T_2=6,\ T_3=2,\ T_4=22,\ T_5=1$，平均 $\bar T=\frac{45}{5}=9$
	\item 响应：$R_3=0,\ R_5=0,\ R_2=3,\ R_1=6,\ R_4=12$，平均 $\bar R=\frac{21}{5}=4.2$
\end{itemize}

\subsection*{(4) MLFQ}

调度序列：

\medskip
$
\begin{aligned}
&1,2,3,4,5,1,2,3,\\
&4,1,2,4,1,4,1,4,\\
&1,4,1,4,1,4,4,4
\end{aligned}
$

\medskip
完成时刻：$C_5=5, C_3=8, C_2=11, C_1=21, C_4=24$

\begin{itemize}
	\item 周转：$T_1=21,\ T_2=11,\ T_3=8,\ T_4=22,\ T_5=3$，平均 $\bar T=\frac{65}{5}=13$
	\item 响应：$R_1=0,\ R_2=1,\ R_3=2,\ R_4=1,\ R_5=2$，平均 $\bar R=\frac{6}{5}=1.2$
\end{itemize}

\medskip

\subsection*{四种算法对比表}
\begin{center}
\begin{tabular}{|c|c|c|}
\hline
调度算法 & 平均周转时间 $\bar T$ & 平均响应时间 $\bar R$ \\
\hline
FCFS & 15 & 10.2 \\
\hline
RR & 13 & 2 \\
\hline
SJF & 9 & 4.2 \\
\hline
MLFQ & 13 & 1.2 \\
\hline
\end{tabular}
\end{center}

\end{document}